\section{Discussion and conclusions}

This study advances the literature on BSS planning by demonstrating how a flexible optimization framework can address multiple planning goals, explicitly manage the spatial arrangement of stations, and integrate topographic and expansion considerations within a single framework. Its versatility was illustrated using a MCDM approach applied to three scenarios targeting demand, first–last kilometer connectivity, and social equity. Building on these scenario-specific outcomes, the integration of station proximity and system accessibility metrics, derived from directed network distances and balanced through a tunable parameter, provided explicit control over the spatial arrangement of stations while still reflecting the underlying planning goals set by the local factor weights. Moreover, incorporating topography-adjusted distances further refined planning in steep districts such as Horta-Guinardó and Nou Barris, enabling a redistribution of stations toward areas that would otherwise face reduced accessibility. Finally, the l’Eixample district expansion experiment demonstrated the framework’s ability to guide incremental growth, ensuring that new stations integrate coherently with existing networks while respecting spatial constraints.

A first aspect concerns the definition of candidate sites. Many studies have defined potential stations at the level of zones or regular grids~\cite{Frade2015, Mete2018, Fazio2021, Veillette2018, Tera2023}, which offers tractable formulations but limits the precision needed for final deployment. Other contributions have moved toward higher resolution by considering all street intersections as candidate sites to better capture local conditions~\cite{Cintrano2018, Xin2013}. We follow this line of work, since planning decisions are ultimately made at the street level, and evaluating the full set of intersections ensures comprehensive coverage of the urban network. In practice, this assumption should be interpreted flexibly: if a node is identified as promising, the actual station would likely be placed in its surroundings where conditions are suitable, such as adjacent sidewalks or plazas. At the same time, this approach excludes mid-block locations and overlooks practical considerations such as the availability of suitable public space for station installation. Future extensions could enrich the computational framework by integrating contextual information—potentially derived from satellite or street-level imagery—to better approximate where stations could realistically be deployed in the urban layout. For example, recent studies have demonstrated how deep-learning models applied to street-view imagery and point clouds can automatically identify and classify street-level elements such as vegetation, roads, sidewalks and street furniture~\cite{Ning2022, Zhou2022}, offering a promising direction for detecting suitable public space for station installation.

Beyond the choice of candidate locations, another key dimension is how planning objectives are framed. Multi-criteria methods have proven valuable for assembling diverse spatial factors such as population and transport accessibility into composite suitability scores~\cite{Kabak2018, Fazio2021, Veillette2018, Tera2023}, but they usually reduce all factors into a single aggregated objective, most often demand coverage. Multi-objective formulations represent an important step forward because they separate competing priorities such as social equity, demand, or multimodal integration and make the trade-offs explicit~\cite{Conrow2018, Duran-Rodas2021, Caggiani2020, Qian2022, Fan2024, Nikiforiadis2021}. Yet, these models also typically rely on a fixed set of goals and spatial factors, so exploring alternative priorities often requires new frameworks or reformulations. Our contribution is not to replace these approaches but to highlight that the multi-criteria logic can be applied more openly: by flexibly deciding the weights and incorporating any data source relevant to the case at hand, planners can generate multi-objective scenarios with a wide range of planning goals within a single framework, avoiding the need to use multiple models. Still, the process of assigning weights remains partly subjective unless supported by stakeholder input or empirical evidence. To address this, decision-making techniques such as AHP, TOPSIS, or MOORA could be incorporated to make weighting more systematic and transparent, as in previous works~\cite{Kabak2018, Guler2021a, Bahadori2022}.

Complementing these aspects, a third contribution concerns the spatial arrangement of stations, a dimension often overlooked in earlier optimization studies. Previous optimization models have typically treated the spatial arrengement of stations as an incidental outcome of the optimization process~\cite{Frade2015, Cintrano2018, Mete2018, Lin2011}. Our framework addresses this by incorporating proximity and accessibility metrics that make spatial arrangement explicit, enabling planners to balance clustered versus dispersed layouts while still reflecting the local factor weights. It relies on a directed graph rather than the undirected representations common in earlier studies~\cite{Cintrano2018, Xin2013}. This is because circulation rules of directed networks such as one-way streets and turn restrictions create asymmetric accessibility, meaning that shortest-path distances can differ depending on travel direction~\cite{Melo2022}. Slope is also accounted for, as it strongly influences cycling behaviour, generally discouraging uphill trips while encouraging downhill ones~\cite{MateoBabiano2016, Grau-Escolano2024}. To our knowledge, García-Palomares et al.~\cite{GarciaPalomares2012} is the only optimization study to incorporate slope, applying a uniform penalty on travel times of climbs. While this was an important advance, the uniform treatment overlooked the different effects of climbs and descents. Our framework refines this with an equivalent flat distance formulation based on the slope–speed relation proposed by Parkin \& Rotheram~\cite{Parkin2010}, penalizing climbs while recognizing that moderate declines can facilitate cycling and steep descents may reduce effective speed. Incorporating directionality and topography together provides a more realistic basis for evaluating inter-station proximity and system accessibility, and thus for designing station networks that better reflect actual cycling conditions.

In addition to modelling choices, the algorithmic approach also deserves reflection. Most \ac{bss} optimization studies that use metaheuristics rely on genetic algorithms, with only a few exploring alternatives such as \ac{simann} or \ac{pso}~\cite{Cintrano2018, Qian2022, Liu2015, Caggiani2020}. While these methods are effective for handling complex objectives, they are typically applied without formal validation against exact formulations, leaving open the question of how close the solutions are to the true optimum. In our case, the GA was adapted to the specific structure of the problem and benchmarked against a linear MILP proxy. The algorithm’s hyperparameters were tuned on a single optimization scenario, yet the comparison performed across multiple settings—combining different spatial factor weights and station counts—showed that the GA consistently achieved very small accuracy gaps relative to the MILP solutions. These results demonstrate that the algorithm generalizes well across planning scenarios while maintaining near-optimal performance when evaluated on the linear component of the objective function. Although this comparison is limited to the linear formulation, it nonetheless provides strong evidence of the GA’s robustness and its capacity to generalise when exact optimisation of the full nonlinear formulation is infeasible. A natural next step for strengthening this validation would be to compare the GA outcomes with the full nonlinear objective function against those from other metaheuristics, thereby testing its robustness more broadly. Beyond validation, further advances in BSS planning could come from the use of multi-objective optimization algorithms specifically designed to approximate Pareto fronts, such as NSGA-II or SPEA2. Unlike weighted-sum approaches, which require multiple runs to explore trade-offs, these algorithms generate sets of non-dominated solutions in a single run, offering a more comprehensive and transparent view of the trade-offs between competing objectives such as demand, social equity, and multimodal integration.

A final consideration concerns the assumptions of the framework. Its effectiveness depends on the availability and quality of open data, which may vary across contexts and affect transferability. Another important point is that urban spatial factors are often correlated, blurring the distinction between separate planning goals. For this reason, preliminary analyses should be conducted to understand these interdependencies and remain aware of the trade-offs that arise when combining different factors into planning objectives.

Taken together, the methodological elements of the framework have direct implications for practice. By flexibly adjusting weights and objectives, planners can use it as a multi-objective scenario-testing tool to explore how different planning priorities affect the spatial distribution of stations, by making trade-offs more explicit. The framework is designed to be transferable because it builds on open datasets such as OSM, census, and transport network data, and a modular structure that allows any spatial factor to be incorporated, enabling its application in other cities with similar data coverage. In addition, although not applied in this study, the framework allows the optimization to focus on specific groups of POIs when planners wish to emphasize particular trip generators, such as hospitals, universities, or business districts. This is made possible by the classification of POIs into categories—such as healthcare, education, culture, or retail—described in the data section. Likewise, proximity to cycling infrastructure can be prioritized by assigning higher weights to the bike-lane variable or by focusing the optimization solely on this factor.

Looking ahead, the framework could be enriched with additional layers of realism to further align with planning practice. One promising direction involves incorporating demand patterns, using data derived from mobility surveys or usage records of an existing BSS. Integrating such information would make it possible to guide the optimization toward areas with demonstrated demand and even to validate alternative scenarios through mobility simulations. In addition, considering the availability of public space and operational constraints such as station capacity would help ensure that proposed sites are not only theoretically optimal but also practically feasible. Political and institutional factors, including budgetary limits linked to station number and size, rebalancing operations, or stakeholder preferences regarding social equity, coverage, or multimodal integration, could likewise be incorporated to make the tool more directly applicable in real-world contexts. Acknowledging these potential extensions underscores the framework’s capacity to evolve into a decision-support system that bridges technical optimization with the complexities of urban governance.